\documentclass[12pt]{article}

\usepackage[utf8]{inputenc}
\usepackage[T1]{fontenc}
\usepackage[polish,provide=*]{babel}
\usepackage{lmodern}
\usepackage{amsmath}
\usepackage{latexsym,amsfonts,amssymb,amsthm,amsmath}
\usepackage{enumitem}
\usepackage{hyperref}
\usepackage{float}
\usepackage[none]{hyphenat}
\usepackage{braket}
\usepackage{graphicx}
\usepackage{subcaption}
\usepackage{booktabs}
\usepackage{physics}
\graphicspath{{./images/}}

\setlength{\parindent}{0in}
\setlength{\oddsidemargin}{0in}
\setlength{\textwidth}{6.5in}
\setlength{\textheight}{8.8in}
\setlength{\topmargin}{0in}
\setlength{\headheight}{18pt}

\title{Superconducting qubit readout: an open quantum system perspective}
\author{Kacper Kłos}

\begin{document}
\maketitle


We start with system Hamiltonian for the transmon interacting with the resonator is (Formula (31) in Blase, ignoring the coupled transmition line) 
\begin{equation}
    \hat{H} = 4 E_C (\hat{n} + \hat{n}_r)^2 - E_J \cos \varphi
\end{equation}
Because $\hat{n} = \frac{\hat{Q}}{2e}$, $\hat{\varphi} = \frac{2e}{\hbar} \hat{\Phi}$, and $[\hat{\Phi}, \hat{Q}] = i$  that means 
\begin{equation}
    [\hat{n}, \hat{\varphi}] = \frac{1}{\hbar} [\hat{Q}, \hat{\Phi}] = -i
\end{equation}
This allow us to see how $e^{i \hat{\varphi}}$ works on state $\ket n$ to compute the $\cos \hat{\varphi} \ket n$. We start with Cambpell lemma:
\begin{equation}
    e^{i \hat{\varphi}} \hat{n} e^{-i \hat{\varphi}} = \hat{n} + [i \hat{\varphi}, \hat{n}] + \frac{1}{2} [i \hat{\varphi}, [i \hat{\varphi}, \hat{n}]] = \hat{n} - 1
\end{equation}
This leads to:
\begin{equation}
    \hat{n}(e^{-i \hat{\varphi}} \ket n) = (n-1) e^{-i \hat{\varphi}} \ket n
\end{equation}
So, the state $(e^{i \hat{\varphi}} \ket n)$ is the eigenstate with the eigenvalue $n+2$. This leads to $(e^{\pm i \hat{\varphi}} \ket n) = \ket{n \pm 1}$. This is relevent for the the $\cos \varphi = \frac{1}{2} (e^{i \hat{\varphi}} + e^{-i \hat{\varphi}})$. Thus, we know how the $H_t$ that is $H$ when $\hat{n}_r = 0$ work on the state $\ket n$, allowing us to describe the matrix.
\begin{equation}
\hat{H_t} \ket n = 4 E_C n^2 \ket n - \frac{E_J}{2}(\ket{n+1} + \ket{n-1})
\end{equation}

\end{document}
